
\subsection{光栅衍射图样，缺级现象}

\begin{theorem}[相邻两个主极大之间有什么？]
    若光栅有 \(N\) 条狭缝：\(N\) 条缝发出的 \(N\) 个等长振动相量首尾相连能够连成一个正多边形，由于这个正多边形的外角为$\Delta\phi$，外角和为$2\pi$，所以一定\[N\Delta\varphi=2\pi q\Leftrightarrow \Delta\varphi=\frac{2\pi q}{N}\]但 \(q=0\) 和 \(q=N\) 对应的是$\Delta\varphi=0,\quad 2\pi$，是主极大，不是暗纹。于是只剩\[\Delta\varphi=\frac{2\pi}{N},\frac{4\pi}{N},\frac{6\pi}{N},\cdots,\frac{2(N-1)\pi}{N}\]共 \(N-1\) 个暗纹。因此相邻两主极大明纹之间有 \(N-1\) 条暗纹； 
真正夹在两条相邻暗纹之间的区间是
\[
\left(\frac{2\pi}{N},\frac{4\pi}{N}\right),
\left(\frac{4\pi}{N},\frac{6\pi}{N}\right),
\cdots,
\left(\frac{2(N-2)\pi}{N},\frac{2(N-1)\pi}{N}\right).
\]这样的区间共有$(N-1)-1=N-2$个。每一个区间两端都是暗纹，光强为零；但区间内部一般不为零，因此光强必然先从零增大，再降回零，中间形成一个局部极大。这个局部极大就是次级明纹。
所以，相邻两主极大明纹之间有$N-2$条次级明纹。
\end{theorem}

\begin{figure}[htbp]
\centering
\includegraphics[width=0.6\textwidth]{pdf-exercises/15.png}
\caption{光栅衍射图样}
\end{figure}
\begin{enumerate}
\item 第一张图“只有衍射”指的是：暂时不考虑许多缝之间的干涉，只看每一条狭缝自身的单缝衍射效果。
单缝衍射的光强分布是
\[
I=I_0\left(\frac{\sin\alpha}{\alpha}\right)^2,
\qquad
\alpha=\frac{\pi a\sin\theta}{\lambda}.
\]中央方向 \(\theta=0\) 最亮；两侧逐渐变暗；
在满足\[
a\sin\theta=\pm k\lambda,\qquad k=1,2,3,\cdots
\]的方向出现暗纹；
中央明纹最宽，两侧次级明纹越来越弱。
\item “只有干涉”指的是：把每条缝理想化成很窄的线光源，不考虑每条缝自身的宽度 \(a\)，只考虑不同缝之间的相干叠加。
当
\[
d\sin\theta=k\lambda,\qquad k=0,\pm1,\pm2,\cdots
\]时，相邻两缝发出的光相位相同，所有缝的振动同向叠加，于是出现主极大明纹。如果有 \(N\) 条缝，那么主极大方向上$I_P=N^2I_0$，但是在两个相邻主极大之间，相邻缝的相位差逐渐变化。若有 \(N\) 条缝，则相邻两个主极大之间会有$N-1$条暗纹，以及$N-2$条次级明纹。
\item 真实光栅不是无限细的线光源，而是由许多有宽度 \(a\) 的狭缝组成。因此真实光栅同时具有两种效应：每一条缝自己会产生单缝衍射包络：
\[
\left(\frac{\sin\alpha}{\alpha}\right)^2.
\]它决定整体强弱分布。许多条缝之间互相干涉，产生尖锐主极大：
\[
\left(\frac{\sin N\beta}{\sin\beta}\right)^2,
\qquad
\beta=\frac{\pi d\sin\theta}{\lambda}.
\]它决定一根根尖锐条纹的位置。
所以真实光强可以写成：
\[
I(\theta)
=
I_0
\left(\frac{\sin\alpha}{\alpha}\right)^2
\left(\frac{\sin N\beta}{\sin\beta}\right)^2.
\]
\end{enumerate}


\begin{definition}[缺级现象]
    光栅方程为$(a+b)\sin\theta=d\sin\theta=\pm k\lambda,k=0,1,2,3,\cdots$，这是屏上可观察级次对应的主极大条件。单缝衍射暗纹条件为$a\sin\theta=\pm k'\lambda,k'=1,2,3,\cdots$，注意这里 \(k'\) 从 \(1\) 开始，因为单缝暗纹没有 \(k'=0\) 这个暗纹，中央方向是中央明纹。如果某个方向同时满足这两个式子，那么把它们相除：
\[
\frac{(a+b)\sin\theta}{a\sin\theta}
=
\frac{k\lambda}{k'\lambda}\implies
\frac{a+b}{a}=\frac{k}{k'}\implies
k=k'\frac{a+b}{a}
\]当某一级 \(k\) 同时等于某个正整数 \(k'\) 乘上 \((a+b)/a\) 时，这一级主极大会缺失。举个很常见的判断：
如果$\frac{a+b}{a}=3$那么$k=3k'$，第 \(3,6,9,\cdots\) 级主极大会缺级。
\end{definition}

\begin{theorem}[屏上可观察的最高级次]
因为真实衍射角不能达到或超过 \(90^\circ\)。由光栅方程 \(d\sin\theta=k\lambda\)，又因为 \(\sin\theta<1\)，所以 \(k\lambda<d\)，即 \(k_{\max}<\frac{a+b}{\lambda}\)。图中写成 \(k_{\max}<\frac{a+b}{\lambda}\)。这表示光栅能看到的最高级次是有限的。不是 \(k=0,1,2,\cdots\) 可以无限往上排，而是受 \(\sin\theta\le 1\) 限制。
实际做题时常写：
\(|k|\le \frac{d}{\lambda}\)，然后取不超过它的最大整数。但如果题目严格写 \(\theta<90^\circ\)，就注意不能取等号对应的极限方向。
\end{theorem}
注意区分“最多能看到几条明纹”和“可观察最高级次”的区别。
\begin{exercise}[500 条每毫米的光栅为什么最后只看见 5 条明纹]
用波长 \(\lambda=\SI{500}{nm}\) 的单色光垂直照射到每毫米有 500 条刻痕的光栅上。求：1. 第一级和第三级明纹的衍射角；2. 若缝宽与缝间距相等，用此光栅最多能看到几条明纹。
\end{exercise}

\begin{solution}
（1）
 \[\sin\theta_1=\frac{\lambda}{d}=0.25,\quad \theta_1\approx14^\circ 28';\qquad \sin\theta_3=\frac{3\lambda}{d}=0.75,\quad \theta_3\approx48^\circ 35'.\]
 （2）极限条件\[d\sin\theta=k\lambda< d\implies k<\frac{d}{\lambda}=4\]
同时题目给 \(a=b\)，即 \(d=a+b=2a\)，缺级条件
 \[d\sin\theta=k\lambda,\qquad a\sin\theta=k'\lambda \implies \frac{d}{a}=\frac{k}{k'}=2\]
说明偶数级主极大都会被单缝暗纹压掉。这样 \(\pm2\) 也不出现，最后只剩 \(k=0,\ \pm1,\ \pm3\)，共 5 条明纹。
\end{solution}

\begin{exercise}[2006级期末：给定光栅常数求可观察最高级次]
波长为 \(\lambda=\SI{550}{nm}\) 的单色光垂直入射于光栅常数 \(d=2\times10^{-4}\,\si{cm}\) 的平面衍射光栅上。求可能观察到光谱线的最高级次。
\end{exercise}

\begin{solution}
光栅主极大由 \(d\sin\theta=k\lambda\) 决定。最高级次要取满足 \(\sin\theta\le1\) 的最大整数 \(k\)。
 \(d=2\times10^{-4}\,\si{cm}=2.0\times10^{-6}\,\si{m}, \qquad \lambda=550\,\si{nm}=5.50\times10^{-7}\,\si{m},\) 
因此
 \(d\sin\theta=k\lambda,\sin\theta\le1 \implies k\le\frac{d}{\lambda} =\frac{2.0\times10^{-6}}{5.50\times10^{-7}}\approx3.64.\) 
可取的最大整数级次为 \(3\)。
\end{solution}

\begin{exercise}[2006级期末：单缝中央明纹宽度内有几个光栅主极大]
一衍射光栅每厘米有 \(200\) 条透光缝，每条透光缝宽 \(a=2\times10^{-3}\,\si{cm}\)。在光栅后放一焦距 \(f=\SI{1}{m}\) 的凸透镜，以 \(\lambda=\SI{600}{nm}\) 的单色平行光垂直照射光栅。求：1. 单缝衍射中央明条纹宽度；2. 在该宽度内有几个光栅衍射主极大。
\end{exercise}

\begin{solution}
这题同时有单缝衍射和光栅主极大。中央明纹宽度先由单缝第一暗纹定边界；在这个角范围内，再用光栅方程数主极大。
\[
a=2\times10^{-3}\,\si{cm}=2.0\times10^{-5}\,\si{m},\qquad
a\sin\varphi=\lambda,\qquad
x\approx f\sin\varphi
\]
所以中央明纹半宽 \(x=f\lambda/a=\SI{0.03}{m}\)，全宽为 \(\Delta x=\SI{0.06}{m}\)。
光栅常数为
\[
d=\frac{1}{200}\,\si{cm}=5.0\times10^{-3}\,\si{cm}=5.0\times10^{-5}\,\si{m}.
\]
要落在单缝中央明纹内，光栅主极大还要满足
 \[\begin{cases}d\sin\theta=k\lambda\\|a\sin\theta|<\lambda\end{cases} \implies \abs{k}\frac{\lambda}{d}<\frac{\lambda}{a} \implies \abs{k}<\frac{d}{a}=2.5.\]
由于$k<\frac{d}{a}$故此时考虑缺级现象不会再删除待选值。因此中央明纹宽度内有 \(k=0,\ \pm1,\ \pm2\) 共 5 个光栅主极大。
\end{solution}

\begin{exercise}
用波长 \(\lambda=\SI{500}{nm}\) 的单色光垂直照射某光栅，观察到第二级主极大出现在 \(\sin\theta=0.20\) 处，且第四级缺级。求：1. 光栅常数；2. 狭缝最小宽度；3. 可见条纹级次。
\end{exercise}

\begin{solution}
\[d\sin\theta=2\lambda\implies d=5\times10^{-6}\text{m}\]
第$k=\frac{d}{a}k'$会缺级，代入$k=4$，有$4=\frac{d}{a}k'$，狭缝宽度$a$最小时，$k'$要尽量小，此时可取最小值$k'=1$，于是
 \[a_{\min}=\frac{d}{4}=\SI{1.25e-6}{m},\qquad \frac{d}{a}=4.\] 
所有满足 \(k=4k'\) 的级次都缺掉，即 \(\pm4,\pm8,\dots\)。同时最高级次还要考虑 \(|\sin\theta|<1\) 限制：
 \[k_{\max}<\frac{d}{\lambda} =\frac{5.0\times10^{-6}}{5.0\times10^{-7}}=10, \qquad k=0,\ \pm1,\ \pm2,\ \pm3,\ \pm5,\ \pm6,\ \pm7,\ \pm9.\]
\end{solution}


\begin{exercise}[两种波长第二次重合时怎样反推光栅常数]
一束平行光垂直入射到某光栅上，其中含有两种波长
 \(\lambda_1=\SI{440}{nm},\qquad \lambda_2=\SI{660}{nm}.\) 
实验发现，两种波长的谱线（不计中央明纹）第二次重合于衍射角 \(\varphi=60^\circ\) 的方向上。求该光栅的光栅常数 \(d\)。
\end{exercise}

\begin{solution}
\textbf{第二次}重合要先找满足
 \(m_1\lambda_1=m_2\lambda_2\) 
的第二组正整数。
 \[440m_1=660m_2 \implies 2m_1=3m_2, \qquad (m_1,m_2)=(3,2),(6,4),\dots\]
不计中央明纹时，第一次重合是 \((3,2)\)，第二次重合是 \((6,4)\)。在 \(\varphi=60^\circ\) 方向上
\[
d\sin 60^\circ=6\lambda_1=2.64\times10^{-6}\ \si{m}.~d=\frac{2.64\times10^{-6}}{\sin 60^\circ} \approx \SI{3.05e-6}{m}
\]
\end{solution}

\begin{exercise}[2015-2016期末：单缝暗纹与光栅主极大怎样分开算]
在单缝夫琅禾费衍射实验中，垂直入射光含有两种波长
 \(\lambda_1=\SI{400}{nm},\lambda_2=\SI{600}{nm}.\) 
已知单缝宽度 \(a=1.0\times10^{-2}\,\si{cm}\)，透镜焦距 \(f=\SI{50}{cm}\)。求两种光第一级衍射暗纹中心之间的距离，可作小角近似。
若用光栅常数 \(d=1.0\times10^{-3}\,\si{cm}\) 的光栅替换单缝，其他条件不变，求两种光第一级主极大之间的距离。
\end{exercise}

\begin{solution}
利用$a\sin\theta=a\frac{x}{f}=\lambda$
 \[x_{\text{单缝}} =f\frac{\lambda_2-\lambda_1}{a} =0.50\times\frac{2.0\times10^{-7}}{1.0\times10^{-4}} =1.0\times10^{-3}\,\si{m} =\SI{1.0}{mm}\]
光栅第一级主极大同理，只是把 \(a\) 换成光栅常数 \(d\)：
 \[\Delta x_{\text{光栅}} =f\frac{\lambda_2-\lambda_1}{d} =0.50\times\frac{2.0\times10^{-7}}{1.0\times10^{-5}} =1.0\times10^{-2}\,\si{m} =\SI{1.0}{cm}\]
\end{solution}


\begin{exercise}[单缝边缘光线相位差为 \(\pi\) 时焦平面位置]
波长 \(\lambda=\SI{480.0}{nm}\) 的平行光垂直照射宽度 \(a=\SI{0.40}{mm}\) 的单缝，单缝后透镜焦距为 \(f=\SI{60}{cm}\)，观察屏位于透镜焦平面上，\(O\) 为透镜焦点，\(P\) 为焦平面上一点。若单缝两边缘点 \(A,B\) 射向 \(P\) 点的两条光线在 \(P\) 点相位差为 \(\pi\)，求距离 \(OP\)。

\begin{center}
\includegraphics[width=0.28\linewidth]{pdf-exercises/pdf2004-q19-single-slit.png}
\end{center}
\end{exercise}

\begin{solution}
两边缘光线相位差 \(\pi\) 表示光程差为 \(\lambda/2\)。所以$a\sin\theta=\frac{\lambda}{2}$
\[\sin\theta=\frac{\lambda}{2a} =\frac{480.0\times10^{-9}}{2\times0.40\times10^{-3}} =6.0\times10^{-4}.\]
 \[x\approx f\sin\theta =0.60\times6.0\times10^{-4}\ \si{m} =3.6\times10^{-4}\ \si{m} =\SI{0.36}{mm}.\]
\end{solution}

\begin{exercise}[钠光与另一光源在同一光栅上怎样比较波长]
用钠光 \(\lambda_1=\SI{589.3}{nm}\) 垂直照射某光栅，测得第三级光谱衍射角为 \(60^\circ\)。若换用另一光源，测得其第二级光谱衍射角为 \(30^\circ\)，求后一光源的波长；再求白光 \(400\sim760\,\si{nm}\) 的第二级光谱张角。
\end{exercise}

\begin{solution}
题目给的是“同一光栅、不同波长、不同级次”的两组数据，因此要用同一个光栅方程把它们联起来。
\[
\begin{aligned}
d\sin\theta&=m\lambda,\qquad
d\sin60^\circ=3\lambda_1,\qquad
d\sin30^\circ=2\lambda_2,\\
\lambda_2&=\frac{3\lambda_1\sin30^\circ}{2\sin60^\circ}
\approx \frac{3\times589.3\times0.5}{2\times0.866}\,\si{nm}
\approx \SI{510.3}{nm},\qquad
d=\frac{2\lambda_2}{\sin30^\circ}\approx \SI{2041.2}{nm}.
\end{aligned}
\]
对白光第二级边界，继续沿用同一个 \(d\)：
\[
\begin{aligned}
\sin\theta_{\min}&=\frac{2\times400}{2041.2}\approx0.392,\qquad
\sin\theta_{\max}=\frac{2\times760}{2041.2}\approx0.745,\\
\theta_{\min}&\approx23^\circ,\qquad
\theta_{\max}\approx48^\circ,
\Delta\theta\approx25^\circ.
\end{aligned}
\]

所以后一光源波长约为 \(\SI{510.3}{nm}\)，白光第二级光谱张角约为 \(25^\circ\)。第二问仍然是同一块光栅，必须继续沿用第一问求出的 \(d\)。
\end{solution}
